\subsection{Horizontal Translation}
To translate a point to the right by $c$, we add $c$ to its $x$-coordinate. To translate it to the left by $c$, we subtract $c$ from its $x$-coordinate.
\begin{Example}
In the figure below, the graph of:
\begin{itemize}
\item $h(x)$ is obtained by translating the graph of $f(x)$ three units to the right
\item $k(x)$ is obtained by translating the graph of $f(x)$ four units to the left
\end{itemize}
Give expressions for $h(x)$ and $k(x)$ in terms of $f(x)$.
\begin{icpic}[x=0.5\linespace,y=0.5\linespace]
\useasboundingbox (-9,-5) rectangle (9,5);
\setgraph{9}{9}{5}{5}
\GraphSmall
\foreach \x/\col in {0/colf,-4/colh,3/colg} {
	\draw[thick,\col] (-4+\x,1) -- (-2+\x,2) -- (2+\x,-4) -- (5+\x,-2);
}
\draw[colf] (-2,2) node[anchor=south,font=\small] {$y=f(x)$};
\draw[colh] (-7,1.5) node[anchor=south east,font=\small] {$y=k(x)$};
\draw[colg] (1.75,1) node[anchor=south west,font=\small] {$y=h(x)$};
\end{icpic}
\funtable{4}{f}{h}{k}
\end{Example}
%%%%%%%%%%%%%%%%%%%%%%%%%%%%%%%%%%%%%%%%%%%%%%%%%%%%%%%%%%
\begin{displaybox}[Horizontal Translation]
For any number $c>0$:
\begin{itemize}
\item If $g(x)=f(x-c)$, the graph of $g$ is the same as the graph of $f$, translated to the \makeblank[1in] by $c$ units.
\item If $g(x)=f(x+c)$, the graph of $g$ is the same as the graph of $f$, translated to the \makeblank[1in] by $c$ units.
\end{itemize}
\end{displaybox}
\begin{Note}
The direction of translation is the \makeblank[1.5in] of the sign of the number.
\end{Note}